\begin{lemma}[Rational projective mixed tensors]
\label{c:super-projective-tensors}
\status{Proved}
Let $A=\Q^{a|b}$ be a finite supervector space, let $L$ be an even
line, and put $G=\operatorname{GL}(A)\times\mathbb G_m$.
We use even morphisms and the category $\mathcal R$ of rational
$G$-supermodules on which
$\epsilon=(\operatorname{diag}(1_a,-1_b),1)$ acts as intrinsic parity.
Every finite object $E$ of $\mathcal R$ is a quotient of a finite
projective-injective object $P$ of $\mathcal R$.  Such $P$, its dual,
and a fixed nonzero projective object $P_0$ are direct summands of
finite sums of mixed tensor powers of $A,A^\vee,L,L^\vee$.
The modules, quotient maps, mixed tensor embeddings and retractions
can be constructed over $\Q$ by terminating finite procedures.
Projectivity also applies to surjections of locally finite rational
modules, including the polynomial modules used below.
\end{lemma}

\begin{proof}
Write $G_0=\operatorname{GL}_a\times\operatorname{GL}_b\times
\mathbb G_m$, $\mathfrak g=\operatorname{Lie}G$, and
$B=U(\mathfrak g_0)$.  We take
\[
 P=U(\mathfrak g)\otimes_B E|_{G_0},\qquad
 P\longrightarrow E,\quad v\otimes e\longmapsto ve.
\]
The odd part of $\mathfrak g$ is
$\operatorname{Hom}(\Q^a,\Q^b)\oplus
\operatorname{Hom}(\Q^b,\Q^a)$.
The PBW theorem makes $P$ finite, of total dimension
$2^{2ab}\dim E$.  Its natural rational $G_0$-action, together with
the compatible $\mathfrak g$-action, integrates to an algebraic
$G$-action over $\Q$.  Here we use
\citet[Corollary~3.6, p.~9; Proposition~5.4 and its proof,
pp.~23--24; Theorem~5.8 and Corollary~5.10, p.~26]{c_masuokashibata2017super}.
The identity $\operatorname{hy}(G)=U(\mathfrak g)$ in characteristic
zero is stated on p.~22 there.  In the proof of Proposition~5.4,
PBW splits off the finite exterior factor, the rational even action
supplies its coefficient functions, and injectivity of the pairing
with the hyperalgebra proves the coassociativity of the resulting
coaction.  Thus the statement applies over $\Q$, not only after
extension to $\C$.  The induced module still satisfies the
$\epsilon$-parity condition.

Induction is left adjoint to restriction to $G_0$.  The latter has
semisimple finite rational representations, so $P$ is projective.
In this case induction and coinduction coincide, with no character
or parity twist.  We give the reason to check this point over $\Q$:
\[
 \det(\mathfrak g_1)
 =\det(\Q^a)^{-b}\det(\Q^b)^a
  \det(\Q^a)^b\det(\Q^b)^{-a}=1,
 \qquad \dim\mathfrak g_1=2ab.
\]
Choose an ordered odd PBW basis.  Projection to its top exterior
coefficient gives a $B$-bimodule map $U(\mathfrak g)\to B$:
the displayed determinant identity is precisely the vanishing of
the possible even-character twist, and the top degree is even.
Multiplication followed by this map is a nondegenerate relative
pairing.  Complementary exterior monomials pair to $\pm1$, while
the commutator correction terms are triangular in exterior degree.
It identifies $U(\mathfrak g)$ with its $B$-linear dual with the
induction/coinduction bimodule structures.  Tensoring with $E$
gives $\operatorname{Ind}(E)\simeq\operatorname{Coind}(E)$.
Right adjunction and semisimplicity over $G_0$ now make $P$
injective as well.  We may take $P_0=\operatorname{Ind}(\Q)$.

For completeness, the tensor-generation assertion can be checked
without assuming $a\geq b$.  Matrix coefficients and the coaction
embedding express every finite rational superrepresentation as a
subquotient of finite sums of mixed tensors and their parity shifts.
Indeed $\Q[G]$ is generated by the entries of the standard
supermatrix, its inverse, and the two powers of the last torus
coordinate.  Its usual localization by the determinants of the two
even blocks has the same generators: an inverse diagonal block is
recovered from the corresponding block of the full inverse by a
finite geometric series in the odd, hence nilpotent, correction.
Every finite set of coefficient functions therefore comes from
bounded mixed tensor degrees.  To retain the correct category,
the operator consisting of intrinsic parity times $\rho(\epsilon)$
splits the full superrepresentation category into $\mathcal R$ and
its parity shift.  It is $+1$ on the mixed tensors in question and
$-1$ on their parity shifts.  Taking its $+1$ summand shows that an
object of $\mathcal R$ is already a subquotient of the unshifted
mixed tensors.  This also explains the category convention in
\citet[section~2, pp.~3--4]{c_heidersdorf2017mixed};
his Theorem~1.1, p.~2, includes the assertion that every projective
is a mixed tensor.

Write $P=U/V$ with $U\subset T$, where $T$ is a finite sum of mixed
tensors.  Projectivity splits $U\twoheadrightarrow P$, so $P$
embeds in $T$; injectivity retracts that embedding.  The dual is
then a summand of $T^\vee$.  These operations are rational.  For an
explicit terminating search, we enumerate finite mixed tensor sums,
compute their rational intertwiner spaces using the even and odd
Lie generators and the torus weights, find an intertwiner of full
rank, and solve the linear equations for its retraction.  The
existence just proved ensures termination.  All the linear
equations have rational coefficients.

Finally $\operatorname{Hom}_{G}(P_0,-)=
\operatorname{Hom}_{G_0}(\Q,-)$ is exact on locally finite rational
modules.  More generally the same assertion holds for the induced
$P$ with $E$ in place of $\Q$.  Exactness follows by placing
finitely many desired preimages in a finite $G_0$-submodule and
using its semisimplicity.  Alternatively, a finite-dimensional
image in a locally finite $G$-module has preimages in a finite
$G$-submodule, so finite projectivity suffices.  This proves the
last assertion without treating an infinite polynomial module as
a finite representation.
\end{proof}

\begin{theorem}[A finite criterion under the sign conjecture]
\label{c:finite-super-criterion}
\status{Proved under the explicitly stated hypothesis $C^+(X)$}
Let $X/\C$ be smooth projective and connected of dimension $n>0$,
let $o\in\CH^n(X)_\Q$ have degree $1$, and put
\[
 a=\dim H^{\rm even}(X,\Q),\qquad
 b=\dim H^{\rm odd}(X,\Q),\qquad
 N=(a+1)(b+1).
\]
Assume $C^+(X)$: a rational algebraic correspondence on $X$ realizes
the even cohomological projector.  Equivalently, some
$u\in\operatorname{End}(h(X))$ realizes cohomological parity
$(-1)^i$ on $H^i(X,\Q)$.  Neither idempotence of its Chow lift nor
membership in the point-and-diagonal ring is assumed.
There is a finite list of rational marked-Frobenius tensor
expressions $P_1,\ldots,P_s$, constructed from the rational cup
table by the procedure below, such that all the cohomological
relations in the requested rings on every power of $X$ hold in
rational Chow if and only if the following finite conditions hold:
\begin{enumerate}[label=(\roman*)]
\item $c_n(T_X)=(a-b)o$ and $\Delta_*o=o\boxtimes o$;
\item $y_\lambda(h(X)^{\otimes N})=0$, where
      $\lambda=((b+1)^{a+1})$ and $y_\lambda$ is its rectangular
      Young symmetrizer;
\item $P_s(h(X),\mu,1,\operatorname{tr},o,\beta^{-1})=0$ for
      every member of the constructed finite list.
\end{enumerate}
The Young symmetrizer is the row symmetrizer followed by the column
antisymmetrizer, with the tensor symmetry of Chow motives; no
choice of its nonzero normalization affects (ii).
Subject to (i), conditions (ii)--(iii) are finite lists of cycles
in the specified point-and-diagonal rings.  No assertion that
rational equivalence itself is algorithmically decidable is made.
\end{theorem}

\begin{proof}
We spell out the superalgebra construction, because exactness of
the invariant functor is unavailable here.  The hypothesis $C^+$
has exactly the usual meaning in
\citet[section~1, p.~1389]{c_tabuada2018sign}; if $v$ realizes the
even projector, $u=2v-\id$ realizes parity.
Put $A=H^*(X,\Q)$ with its parity, $L=\Q(-1)$ as an even line,
and $G=\operatorname{GL}(a|b)\times\mathbb G_m$.
All tensor representations used in the covariant construction,
including $W$, $R$ and $I$ below, belong to the category $\mathcal R$
of Lemma~\ref{c:super-projective-tensors}.  The entire regular
representation $\mathcal O(G)$ need not belong to that category.
The same marked tensors as in
Theorem~\ref{c:finite-even-criterion} give an even rational point
$t=(\mu,1,\operatorname{tr},o,\beta^{-1})$ in
\[
 W=(A\otimes A^\vee\otimes A^\vee)\oplus A
   \oplus(A^\vee\otimes L^n)
   \oplus(A\otimes L^{-n})
   \oplus(A\otimes A\otimes L^{-n}).
\]
Write $R=\mathcal O(W)$ and define
\[
 I=\ker\bigl(R\longrightarrow\mathcal O(G)\bigr),
 \qquad f\longmapsto\bigl(g\longmapsto f(g\cdot t)\bigr).
\]
Thus $I$ is the ideal of the schematic super-orbit closure.  We
retain its odd directions and nilpotents; the ideal of the orbit
of ordinary complex points would not suffice.

The supercommutative algebra
$R=\operatorname{Sym}(W_{\bar0}^\vee)\otimes
\bigwedge(W_{\bar1}^\vee)$ is finite over its ordinary even
polynomial subalgebra and is noetherian.  We can compute generators
of $I$ by orbit-graph elimination.  Coordinates on $G$ consist of
the two even matrix blocks, the odd off-diagonal blocks, and the
inverses of the two even determinants and of the last torus
coordinate.  Introduce ordinary variables for those inverses and
their defining equations.  Since the odd variables have a finite
exterior basis, elimination is a finite-module polynomial
elimination with explicitly known signed multiplication matrices.
The usual noetherian termination of module Gr\"obner elimination
therefore applies.  Enlarge the resulting finite generating span
to a finite $G$-stable submodule $E\subset I$.  The action preserves
polynomial degree, so this is a finite linear computation; one can
close the span under the rational coaction or the even and odd Lie
generators together with the torus weights.

Use Lemma~\ref{c:super-projective-tensors} to construct a finite
projective-injective $P\twoheadrightarrow E$ and a split mixed
tensor embedding of $P^\vee$.  The composite $P\to E\hookrightarrow R$
corresponds to the polynomial covariant
\[
 p_P\in(P^\vee\otimes R)^G.
\]
Its components in the chosen mixed tensor sum are our finite
expressions $P_1,\ldots,P_s$.  We also fix the nonzero projective
$P_0=\operatorname{Ind}_{G_0}^{G}(\Q)$ and a mixed tensor projector
whose image is $P_0$.

We verify both the diagram interpretation and its exact relations.
\citet[Theorem~2.1, Corollary~2.2, Theorem~2.3 and
Corollary~2.5]{c_lehrerzhang2021super} state the first and second
fundamental theorems for the general linear supergroup in this
scope: mixed invariants are contractions, and their kernel is the
tensor ideal generated by the Young symmetrizer with $a+1$ rows
and $b+1$ columns.  Although the article's title concerns an
orthosymplectic supergroup, these numbered statements are explicitly
its preliminary general linear theorems.  To pass from the
endomorphism version to all mixed tensors we bend dual indices;
the central scalar in $\operatorname{GL}(A)$ forces the two
remaining tensor degrees to be equal.  The stated kernel in
$\Q[\mathfrak S_r]$ is the ideal generated by the rectangular
symmetrizer in the first $N$ tensor positions.  Bending back gives
exactly its tensor ideal in all mixed degrees.  The scalar loop
has value $a-b$.

For polynomial covariants, the supersymmetric power in each finite
degree is the image of the averaging idempotent for the signed
symmetric-group action on the corresponding tensor power of
$W^\vee$.  That idempotent provides a split inclusion before taking
invariants.  Polarization and the preceding tensor theorem
therefore express the covariant as a finite contraction diagram;
no exactness of supergroup invariants is needed.  This is also the
superpolarization proved in Appendix~A, Proposition~A.1 and
Corollary~A.2 of the same source.  The source states these theorems
over $\C$, but all invariant spaces, maps and symmetrizers here
are defined by rational linear equations.  The equality of their
images and kernels descends along $\Q\subset\C$.
The extra torus just enforces the Tate degree.  The fixed Tate
line already satisfies $\bigwedge^2L=0$ and
$\operatorname{tr}(\id_L)=1$.

We now replace the invalid reductive invariant argument.  Let
$F\in(T'\otimes R)^G$ be any polynomial diagram covariant that
vanishes at $t$, where $T'$ is a mixed tensor target.  Equivariance
on superalgebra-valued points implies that its coefficients lie
in $I$, so $F\in(T'\otimes I)^G$.  Multiplication gives an
equivariant surjection $P\otimes R\twoheadrightarrow I$.
Regard $F\otimes\id_{P_0}$ as a $G$-map
\[
 P_0\longrightarrow T'\otimes I\otimes P_0.
\]
Projectivity lifts it through
\[
 T'\otimes P\otimes R\otimes P_0
     \twoheadrightarrow T'\otimes I\otimes P_0.
\]
The lift has bounded polynomial degree: its domain is finite and
the target is locally finite.  The preceding invariant theorem
writes it as a finite diagram between the chosen mixed tensor
summands.  Composing with $p_P$ and multiplying polynomial factors
recovers $F\otimes\id_{P_0}$.  Consequently every relation,
after tensoring with this one fixed $P_0$, lies in the tensor ideal
generated by the finite list $P_s$ and the rectangular universal
relation.  This is a projective lifting argument, not an assertion
that $(\cdot)^G$ is exact.

Evaluate in rational Chow motives, assuming (i)--(iii).  The loop
value is already $\chi(X)=a-b$.  Condition (ii) kills all universal
mixed tensor relations.  In particular, a diagram lift $e$ of the
projector defining $P_0$ is now an actual Chow idempotent: the
universal identity $e^2-e$ is among those relations.  Thus
$P_0(h(X))=\operatorname{im}e$ is defined.  Condition (iii) gives
\[
 F(h(X))\otimes\id_{P_0(h(X))}=0.
\]
It remains to cancel the projective factor, whose categorical
dimension can itself be zero.

Let $u$ be the algebraic correspondence realizing parity supplied
by $C^+(X)$.  On the mixed tensor sum $T(h(X))$ containing $P_0(h(X))$
we tensor $u$ and its dual, and use the identity on the even Tate
line; denote the resulting endomorphism by $u_T$.
Set $u_0=e u_T e$ on $P_0(h(X))$.  Its Betti realization is
parity on $P_0(A)$, because the even projector $e$ commutes with
parity in supervector spaces.  Hence
\[
 \operatorname{tr}(u_0)
  =\operatorname{str}\bigl((-1)^{\rm parity}|P_0(A)\bigr)
  =\dim P_0(A)_{\bar0}+\dim P_0(A)_{\bar1}>0.
\]
The first equality is an equality of rational scalars: Betti
realization is injective on $\operatorname{End}(\mathbf1)=\Q$.
Compose the preceding vanishing with $\id\otimes u_0$ and take
the partial categorical trace.  We obtain
$\operatorname{tr}(u_0)F(h(X))=0$, whence $F(h(X))=0$.
No equality $u^2=1$ in Chow, and no commutation of $u_T$ with $e$
in Chow, was used.  This proves sufficiency for every diagram
relation and therefore for every requested relation.

Conversely, all requested relations imply (i) by
$\Delta\cdot o_1=o_1o_2$ and
$\Delta^2=(a-b)o_1o_2$ and their pushforwards.  The Young
symmetrizer in (ii) is a finite rational combination of permutation
graphs on $X^N$, each a product of the allowed diagonals on
$X^{2N}$, and its Betti class vanishes.  Thus (ii) follows as well.
After (i), the Frobenius graph normalization in the proof of
Theorem~\ref{c:finite-even-criterion} reduces every dualized or
contracted marked diagram to a rational combination of pointed
partitions: each handle inserts $c_n(T_X)$, each marking inserts
$o$, and the two identities in (i) remove these insertions or turn
them into external point classes.  Closed components give their
degrees.  Every $P_s$ vanishes at $t$, so its normalized cycle is
one of the requested cohomological relations.  It therefore
vanishes in Chow.  This proves necessity and the finite iff.
\end{proof}

\begin{remark}[What the additional finite criterion settles]
\label{c:super-criterion-scope}
\status{Proved scope clarification; unrestricted classification remains unresolved}
Theorem~\ref{c:finite-super-criterion} supplies a finite cycle
criterion for mixed cohomology whenever $C^+(X)$ is known.  The
parity correspondence may be completely external to the requested
ring, as it is for the positive pair of
Proposition~\ref{c:product-parity}.  Thus that example's obstruction
to a tautological parity projector does not obstruct this proof.
We do not infer $C^+(X)$ from the point-and-diagonal relations.
Theorem~\ref{c:finite-unconditional-criterion} below removes that
hypothesis from the finite criterion by a different scalar-form
construction, which needs no cancellation step.
Corollary~\ref{c:finite-weil-criterion} removes the rational Betti
restriction for finite existence; its effectiveness qualification
is essential.
The criterion does not prove the finite Chow vanishing conditions
for an arbitrary input variety.  The general classification in
C.1(c), and in particular the K3 vanishing below, remain open
items of this task rather than consequences of finite generation.

Condition (ii) retains a precise finite-dimensionality obstruction.
Under $C^+_H(X)$ it is equivalent to Kimura finite-dimensionality,
by \citet[Theorem~3.2, p.~286]{c_delpadronemazza2005}.
To check the input, their Proposition~3.1 applies the partial-trace
identity of Theorem~2.1 to an algebraic even projector: all its
power traces equal $a$, which lies outside the forbidden contents
for the minimal $(a+1)$-by-$(b+1)$ rectangle.  The numerical ideal
is therefore nilpotent.  The same rectangle identity propagates
to tensor powers by the super tensor relations, so their
numerical ideals are nilpotent as well.  The even and odd
cohomological projectors lift across these nilpotent ideals, and
the cohomologically zero exterior and symmetric projectors on
their tensor powers are nilpotent idempotents and hence zero.
This gives the Kimura decomposition.  Conversely its even and
odd dimensions give the rectangular Schur vanishing.  In
particular, arbitrary Schur-finiteness with an unspecified larger
partition is not being substituted for this exact condition.
\end{remark}
